% =============================================
% CHAPITRE 4 : Implémentation et Évaluation
% =============================================
\chapter{Implémentation et Évaluation}
\label{chap:implementation}

\section{Introduction}

Ce chapitre présente la mise en œuvre du système. Il décrit d'abord l'environnement technique retenu, puis les principaux mécanismes développés, à savoir le gestionnaire de confiance, la sélection des pairs et les mécanismes de sécurité. Les échanges JSON assurant la communication entre les nœuds sont ensuite présentés. Enfin, les résultats des tests et du scénario de mitigation collaborative sont exposés.

% --------------------------------------------------
\section{Environnement technique}
\label{sec:env_tech}



% ==================================================
\subsection{Langages de programmation}
% ==================================================

\techentry{%
  \textbf{JavaScript :} est un langage de programmation interprété, orienté événement et non bloquant. Il constitue le langage principal du back-end via Node.js~\parencite{nodejs2024}, dont le modèle asynchrone est particulièrement adapté aux échanges inter-nœuds concurrents (heartbeats, requêtes d'aide simultanées).%
}{javascript.png}

\techentry{%
 \textbf{SQL :} est le langage de requête standard des bases de données relationnelles. Il est utilisé via PostgreSQL et l'ORM Sequelize pour la définition des schémas, les migrations et l'interrogation des quatorze tables du modèle de données.%
}{sql.png}

% ==================================================
\subsection{Outils de développement}
% ==================================================

\techentry{%
  \textbf{Node.js v20 LTS~\parencite{nodejs2024} :} est une plate-forme côté serveur construite sur le moteur V8 de Chrome. Son modèle événementiel est idéal pour les applications à I/O concurrentes telles que les coalitions de nœuds échangeant des heartbeats sur événement (adhésion, mise à jour, départ) et des requêtes d'aide simultanées.%
}

\techentry{%
 \textbf{Express.js v4~\parencite{expressjs2024} :} est un framework REST minimaliste pour Node.js, de type \textit{middleware-first}. Il structure les routes de l'API (\texttt{/peers}, \texttt{/help}, \texttt{/trust}, \texttt{/auth}) et facilite l'intégration des middlewares mTLS et JWT.%
}

\techentry{%
  \textbf{PostgreSQL 15~\parencite{postgresql2024} :} est un système de gestion de bases de données relationnelles open-source reconnu pour sa conformité ACID. Il est retenu pour son support natif des types \texttt{UUID}, \texttt{ENUM}, \texttt{FLOAT} et la gestion des contraintes de clés étrangères avec politique de suppression en cascade.%
}

\techentry{%
  \textbf{Sequelize v6~\parencite{sequelize2024} :} est un ORM Node.js qui fournit une abstraction SQL avec migrations, associations déclaratives et transactions. Il mappe les quatorze classes du modèle objet vers les tables PostgreSQL.%
}

\techentry{%
 \textbf{jsonwebtoken~\parencite{jsonwebtoken2024} (JWT) :} est la bibliothèque Node.js d'implémentation du standard JWT (RFC~7519 \parencite{rfc7519jwt}). Elle est utilisée pour générer et vérifier les tokens RS256 à durée de vie de 15~minutes qui sécurisent chaque appel API inter-nœuds.%
}

\techentry{%
   \textbf{Swagger / OpenAPI 3~\parencite{openapi3} :} est un standard de description d'API REST. Il auto-documente l'ensemble des endpoints via des annotations JSDoc et expose une interface interactive accessible à \texttt{/api-docs}.%
}

% ==================================================
\subsection{Environnement de déploiement}
% ==================================================

\techentry{%
  \textbf{Docker \& Docker Compose~\parencite{docker2024} :} permettent la containerisation du système. Chaque nœud de la coalition s'exécute dans un conteneur isolé~; Docker Compose orchestre les quatre nœuds et la base de données PostgreSQL avec un réseau virtuel dédié, garantissant un déploiement reproductible sur toute machine disposant de Docker.%
}

\subsection{Architecture de déploiement}

La Figure~\ref{fig:deployment} présente l'architecture de déploiement du prototype. Le système est composé de cinq conteneurs Docker orchestrés par Docker Compose~: quatre nœuds API (\texttt{shieldnet-university}, \texttt{shieldnet-pme}, \texttt{shieldnet-isp},\\ \texttt{shieldnet-datacenter}), chacun exposant son API sur un port distinct (3001 à 3004) en interne sur le port 8443, et une instance PostgreSQL partagée hébergeant quatre bases de données indépendantes. Les nœuds communiquent entre eux via un réseau coalition sécurisé en HTTPS/mTLS (TLS~1.3, authentification JWT RS256) selon trois flux principaux~: le heartbeat événementiel (adhésion, mise à jour de paramètres ou départ — jamais de minuteur périodique), les requêtes d'aide lors d'une attaque, et la collecte de crédibilité réseau pour le calcul du score de confiance.

\begin{figure}[H]
\centering
\figborder[\textwidth]{figures/deployment_diagram1.pdf}
\caption{Diagramme de déploiement du prototype.}
\label{fig:deployment}
\end{figure}

% ==================================================
\subsection{Outils de collaboration}
% ==================================================

\techentry{%
\textbf{Git~\parencite{git2024} :} est le système de contrôle de version distribué utilisé pour gérer l'historique du code source, les branches de fonctionnalités et les correctifs tout au long du développement du prototype.%
}

\techentry{%
\textbf{GitHub~\parencite{github2024} :} est la plateforme de collaboration hébergeant le dépôt du projet. Elle permet le suivi des issues, la revue de code par Pull Requests et la sauvegarde centralisée du code source.%
}

\subsection{Structure du projet}

Le projet est organisé selon le schéma de la Figure~\ref{fig:structure}, séparant la logique métier (\texttt{utils/}), les routes \gls{api} (\texttt{routes/}), les modèles de données (\texttt{models/}) et la configuration (\texttt{config/}).

\begin{figure}[H]
\centering
\figborder[\textwidth]{figures/shieldnet_directory_tree.png}
\caption{Structure des répertoires du projet.}
\label{fig:structure}
\end{figure}


% --------------------------------------------------
\section{Implémentation du gestionnaire de confiance}
\label{sec:impl_trust}

\subsection{Calcul de T(u) avec Cr historique}

Le module \texttt{utils/trustManager.js} implémente la formule PeerTrust (Équation~\ref{eq:trust}). La fonction \texttt{computeTrustScore(peer\_id)} charge toutes les sessions complétées du pair et calcule la moyenne pondérée~:

\begin{lstlisting}[language=JavaScript, caption={Calcul de T(u) avec crédibilité historique par session.}, label=lst:trust]
async function computeTrustScore(peer_id) {
  const sessions = await HelpSession.findAll({
    where: { helping_peer_id: peer_id, status: "COMPLETED" },
    include: [{ model: Attack, as: "attack",
                attributes: ["severity"] }],
  });

  if (sessions.length === 0)
    return { score: DEFAULT_TRUST, level: "BRONZE", session_count: 0 };

  let numerator = 0, denominator = 0;

  for (const session of sessions) {
    const accepted = Number(session.accepted_volume_gbps ?? 0);
    const actual   = Number(session.actual_volume_gbps   ?? 0);

    // S(u,i) = min(1, V_reel / V_accepte)
    const S = accepted > 0 ? Math.min(1, actual / accepted) : 0;

    // D(u,i) : facteur contextuel selon la severite
    const severity = session.attack?.severity ?? "LOW";
    const D = SEVERITY_WEIGHTS[severity] ?? 0.25;

    // Cr_i : credibilite capturee a la cloture de la session
    const Cr_i = session.cr_value ?? DEFAULT_TRUST;

    numerator   += S * Cr_i * D;
    denominator += Cr_i * D;
  }

  const score = denominator > 0 ? numerator / denominator : DEFAULT_TRUST;
  return { score, level: scoreToLevel(score),
           session_count: sessions.length };
}
\end{lstlisting}

La valeur \texttt{cr\_value} est persistée dans \texttt{HELP\_SESSIONS} à la clôture de chaque attaque (route \texttt{POST /attack/over}), comme décrit à la section~\ref{subsec:capture_cr}.

\subsection{Capture de la crédibilité à la clôture}
\label{subsec:capture_cr}

La fonction \texttt{fetchCrFromNetwork(localNodeId, excludePeerId)} interroge tous les pairs actifs (sauf le pair évalué) pour obtenir leur score de confiance courant du nœud local, puis retourne la moyenne~:

\begin{lstlisting}[language=JavaScript, caption={Capture de Cr et persistance à la clôture de session.}, label=lst:cr_capture]
// Dans POST /attack/over -- routes/coalition.js
const localNodeId = await getLocalNodeId();
const crCache = {};  // evite les appels en double par pair

for (const session of sessions) {
  // ... mise a jour actual_volume_gbps ...

  if (localNodeId) {
    const pid = session.helping_peer_id;
    if (crCache[pid] === undefined)
      crCache[pid] = await fetchCrFromNetwork(localNodeId, pid);

    await session.update({ cr_value: crCache[pid] });
  }

  // ... mise a jour du ledger de reciprocite ...
}
\end{lstlisting}

Le cache \texttt{crCache} évite d'interroger le réseau plusieurs fois pour un même pair si plusieurs sessions concernent ce pair dans la même attaque. Le résultat est stocké dans \texttt{HELP\_SESSIONS.cr\_value} (FLOAT, valeur par défaut 0.5 pour les sessions antérieures à cette version).

\subsection{Persistance et sanctions automatiques}

La fonction \texttt{recalculateAndSave(peer\_id)} recalcule le score, le persiste dans \texttt{TRUST\_SCORES} et applique automatiquement la sanction si le score tombe sous 0.20. Le champ \texttt{trust\_level} de la table \texttt{TRUST\_SCORES} est mis à jour~; si le niveau devient BANNED, le pair est marqué \texttt{status = BANNED} et \texttt{membership\_status = EXPELLED} dans la table \texttt{PEERS}. Tout changement de niveau déclenche par ailleurs une entrée dans \texttt{AUDIT\_LOG} (type \texttt{TRUST\_LEVEL\_CHANGED}), et un bannissement automatique génère spécifiquement une entrée de sévérité CRITICAL (\texttt{PEER\_BANNED}).

% --------------------------------------------------
\section{Implémentation de la sélection des pairs}
\label{sec:impl_wsm}

\subsection{Calcul du score WSM}

Le module \texttt{utils/peerSelector.js} implémente l'Équation~\ref{eq:wsm}. La fonction \texttt{selectPeers()} charge par défaut tous les pairs \texttt{ACTIVE} avec capacité disponible, normalise les valeurs brutes, puis calcule le score composite de chaque candidat. Elle accepte également un paramètre optionnel \texttt{peerIds} qui restreint le calcul à un sous-ensemble déjà connu (en sautant le filtrage d'éligibilité) et un paramètre \texttt{capacityOverrides} qui substitue au critère $C_p$ une valeur plus récente que \texttt{declared\_available\_gbps} — c'est ce même mécanisme que réutilise \texttt{POST /attack/\{id\}/allocate} (section~\ref{subsec:allocate}) pour calculer l'allocation sur les seuls pairs ayant accepté, à partir de leur \texttt{accepted\_volume\_gbps}~:

\begin{lstlisting}[language=JavaScript, caption={Calcul du score WSM et normalisation des critères.}, label=lst:wsm]
const W = { C: 0.545, L: 0.193, T: 0.193, R: 0.069 };
const EPSILON = 1e-6;

// Normalisation des valeurs brutes
const maxCap  = Math.max(...capValues);
const minLat  = Math.min(...latValues);
const latRange = Math.max(...latValues) - minLat;

const scored = candidates.map((peer, i) => {
  const Cp     = maxCap    > 0 ? capValues[i] / maxCap               : 0;
  const Lp_inv = latRange  > 0 ? 1 - (latValues[i]-minLat)/latRange  : 1;
  const Tp     = peer.trust_score?.overall_score ?? 0.5;

  const received = Number(peer.reciprocity_ledger?.credits_received ?? 0);
  const given    = Number(peer.reciprocity_ledger?.credits_given    ?? 0);
  const Rp = received / (received + given + EPSILON);

  const score = W.C*Cp + W.L*Lp_inv + W.T*Tp + W.R*Rp;
  return { peer, score, criteria: { Cp, Lp_inv, Tp, Rp } };
});
\end{lstlisting}

\subsection{Plan de répartition proportionnelle}

Le plan de répartition est construit en normalisant les scores WSM~; l'allocation en pourcentage est la partie entière de $w_i \times 100$~:

\begin{lstlisting}[language=JavaScript, caption={Construction du plan de répartition proportionnelle.}, label=lst:alloc]
const totalScore = scored.reduce((sum, c) => sum + c.score, 0);

const plan = scored.map((c) => {
  // w_i = Score(p_i) / sum Score(p_j)
  const weight = totalScore > 0
                 ? c.score / totalScore
                 : 1 / scored.length;

  // allocation_pct = min(floor(w_i * 100), cap_disp)
  const wsm_pct      = Math.floor(weight * 100);
  const allocation_pct = Math.min(wsm_pct, c.cap_disp);

  return { peer: { ... }, wsm_score, weight, allocation_pct };
});
\end{lstlisting}

Le quota de paquets alloué à chaque pair par cycle de Round-Robin pondéré est la partie entière du poids $w_i$ exprimé en pourcentage, plafonné par la capacité disponible du pair~:

\begin{equation}
\label{eq:alloc_cap}
\mathrm{alloc}_i = \min\!\left(\lfloor w_i \times 100 \rfloor,\; \mathrm{cap\_disp}_i\right)
\end{equation}

\noindent Cette contrainte assure qu'un pair ne peut recevoir plus de paquets par cycle que sa capacité déclarée ne le permet, indépendamment de la qualité de son score WSM. La Figure~\ref{fig:wsm_distribution} synthétise l'ensemble du plan de répartition~: depuis le trafic d'attaque entrant jusqu'à l'allocation proportionnelle par pair, en passant par le calcul des poids $w_i$ et les contraintes associées.

% --------------------------------------------------
\subsection{Distribution des paquets par Round-Robin pondéré}
\label{subsec:wrr}

Une fois les allocations $\mathrm{alloc}_i$ calculées pour chaque pair $N_{p_i}$, la redistribution effective des paquets est réalisée selon un mécanisme de \textbf{Round-Robin pondéré} (Weighted Round-Robin, WRR). Le nœud victime $N_v$ dispose d'un flux entrant de $n$ paquets $(P_1, P_2, \ldots, P_n)$ qu'il distribue vers $m$ nœuds pairs $(N_{p_1}, N_{p_2}, \ldots, N_{p_m})$ acceptés, chacun associé à un quota $\mathrm{alloc}_i$ représentant le nombre de paquets qu'il absorbe par cycle.

\paragraph{Principe d'un cycle.}
Les pairs sont parcourus séquentiellement dans l'ordre $N_{p_1}, N_{p_2}, \ldots, N_{p_m}$. À chaque passage, le pair $N_{p_i}$ reçoit exactement $\mathrm{alloc}_i$ paquets consécutifs. Le pair $N_{p_i}$ reçoit ainsi les paquets d'indice~:

\begin{equation}
\label{eq:wrr_range}
P_{\,1+\sum_{j=1}^{i-1}\mathrm{alloc}_j},\quad P_{\,2+\sum_{j=1}^{i-1}\mathrm{alloc}_j},\quad \ldots,\quad P_{\,\sum_{j=1}^{i}\mathrm{alloc}_j}
\end{equation}

Après $N_{p_m}$, le processus revient à $N_{p_1}$ et entame un nouveau cycle. Cette opération est répétée jusqu'à ce que le dernier paquet $P_n$ soit distribué. La Figure~\ref{fig:wrr} illustre ce mécanisme.

\begin{figure}[H]
\centering
\figborder[\textwidth]{figures/wrr_distribution.pdf}
\caption{Distribution des paquets par Round-Robin pondéré~: le nœud victime $N_v$ distribue $n$ paquets vers $m$ pairs en cycles successifs, chaque pair $N_{p_i}$ recevant $\mathrm{alloc}_i$ paquets par cycle.}
\label{fig:wrr}
\end{figure}
% --------------------------------------------------
\section{Mécanismes de sécurité}
\label{sec:impl_secu}

\subsection{JWT RS256~: génération et vérification}

Chaque nœud génère ses tokens avec sa clé privée RSA~; la vérification utilise la clé publique de l'émetteur, stockée dans \texttt{PEERS.public\_key}~:

\begin{lstlisting}[language=JavaScript, caption={Génération et vérification JWT RS256.}, label=lst:jwt]
// Generation (noeud emetteur)
function generateToken(node_id) {
  const privateKey = fs.readFileSync(process.env.TLS_KEY);
  return jwt.sign(
    { iss: node_id, node_id },
    privateKey,
    { algorithm: "RS256", expiresIn: "15m" }
  );
}

// Verification (noeud recepteur) -- extrait du middleware
const decoded   = jwt.decode(token);          // lecture iss sans verif
const peer      = await Peer.findOne({ where: { peer_id: decoded.iss } });
const publicKey = peer.public_key;            // PEM stocke en base
jwt.verify(token, publicKey, { algorithms: ["RS256"] });
\end{lstlisting}

Le token expire après 15 minutes~; un token expiré retourne HTTP 401 avec un message invitant à obtenir un nouveau token via \texttt{POST /auth/token}.

\subsection{mTLS~: configuration côté serveur}

Le serveur HTTPS est créé avec les options TLS suivantes~:

\begin{lstlisting}[language=JavaScript, caption={Configuration mTLS du serveur HTTPS.}, label=lst:mtls]
const server = https.createServer({
  key:  fs.readFileSync(process.env.TLS_KEY),
  cert: fs.readFileSync(process.env.TLS_CERT),
  ca:   fs.readFileSync(process.env.TLS_CA),
  requestCert:          true,  // demander le certificat client
  rejectUnauthorized:   true,  // refuser tout cert non signe par la CA coalition
}, app);
\end{lstlisting}

Lorsque \texttt{MTLS\_ENABLED=true}, \texttt{rejectUnauthorized} est toujours positionné à \texttt{true}~: toute connexion dont le certificat client n'est pas signé par la CA de la coalition est rejetée au niveau TLS avant même d'atteindre les middlewares applicatifs. La variable d'environnement \texttt{MTLS\_STRICT} intervient quant à elle au niveau du middleware applicatif~: si elle vaut \texttt{true}, une requête sans certificat (certificat absent plutôt qu'invalide) est également rejetée avec HTTP 401.

% --------------------------------------------------
\section{Protocole de communication inter-nœuds}
\label{sec:protocole_com}

Les nœuds communiquent exclusivement via des requêtes HTTPS authentifiées par JWT RS256. Cette section présente les messages JSON échangés lors des principales interactions du protocole.

% --------------------------------------------------
\subsection{Authentification}

Avant tout appel API protégé, un nœud obtient un token JWT en transmettant son identifiant et son secret local.

\begin{lstlisting}[language=json, caption={Requête et réponse — \texttt{POST /auth/token}.}, label=lst:auth]
// Requête
{
  "node_id":     "node-university",
  "node_secret": "secret-key-2025"
}

// Réponse 200 OK
{
  "token":      "eyJhbGciOiJSUzI1NiJ9...",
  "expires_in": "15m",
  "node_id":    "node-university",
  "role":       "local",
  "algorithm":  "RS256"
}
\end{lstlisting}

Le token expire après 15 minutes. Le nœud récepteur vérifie la signature RS256 à l'aide de la clé publique de l'émetteur~; un token expiré ou altéré retourne HTTP 401. Le nœud récepteur vérifie la signature RS256 à l'aide de la clé publique de l'émetteur, enregistrée dans la table \texttt{Peers}.

% --------------------------------------------------
\subsection{Heartbeat}

Chaque nœud diffuse un heartbeat à ses pairs connus sur événement uniquement~: adhésion à la coalition (annonce initiale au démarrage, après un court délai), mise à jour de ses propres paramètres (capacité, charge), ou départ propre. Aucun minuteur périodique ne déclenche d'envoi en l'absence de changement~; cette conception évite la croissance en $O(N^2)$ du trafic de supervision observée avec un heartbeat périodique classique lorsque la taille de la coalition augmente.

\begin{lstlisting}[language=json, caption={Requête et réponse — \texttt{POST /heartbeat}.}, label=lst:heartbeat]
// Requête
{
  "peer_id":                  "550e8400-e29b-41d4-a716-446655440000",
  "reported_status":          "ACTIVE",
  "reported_load_pct":        42.5,
  "reported_available_gbps":  5.75,
  "round_trip_time_ms":       12
}

// Réponse 201 Created
{
  "message":   "Heartbeat recorded",
  "heartbeat": {
    "peer_id":                 "550e8400-...",
    "reported_status":         "ACTIVE",
    "reported_load_pct":       42.5,
    "reported_available_gbps": 5.75,
    "round_trip_time_ms":      12
  }
}
\end{lstlisting}

À la réception, le nœud met à jour le champ \texttt{last\_heartbeat} et la capacité déclarée du pair émetteur. Si une attaque est en cours, cette réception déclenche également une reconsidération de la coalition (\texttt{reconsiderCoalition()}, cf.~section~\ref{subsec:reconsideration}), afin de tenir compte immédiatement de la capacité mise à jour. La détection d'un pair injoignable reste réactive~: aucun marquage \texttt{INACTIVE} n'est appliqué sur la seule absence de heartbeat~; il faut l'échec d'une sollicitation réelle (\texttt{POST /help/request}) pour déclencher ce marquage.

% --------------------------------------------------
\subsection{Détection d'une attaque}

Lorsqu'une attaque \gls{ddos} est détectée, le nœud victime l'enregistre localement via \texttt{POST /alert}.

\begin{lstlisting}[language=json, caption={Requête et réponse — \texttt{POST /alert}.}, label=lst:alert]
// Requête
{
  "target_ip_range":  "192.168.1.0/24",
  "target_service":   "HTTP",
  "target_port":      80,
  "target_protocol":  "TCP",
  "status":           "DETECTED",
  "coalition_helped": false
}

// Réponse 201 Created
{
  "attack_id":        "a1b2c3d4-...",
  "status":           "DETECTED",
  "detected_at":      "2026-05-22T10:00:00.000Z",
  "target_ip_range":  "192.168.1.0/24",
  "coalition_helped": false
}
\end{lstlisting}

La sévérité est initialisée à \texttt{LOW} et recalculée à la clôture de l'attaque en fonction du volume total filtré par la coalition.

% --------------------------------------------------
\subsection{Demande et réponse d'aide}

Le nœud victime diffuse une demande d'aide à \emph{tous} les pairs en statut \texttt{ACTIVE} de la coalition. À ce stade, aucune allocation ni volume n'est connu~: la requête se limite à identifier l'attaque et le pair sollicité.

\begin{lstlisting}[language=json, caption={Requête et réponse — \texttt{POST /help/request} (sollicitation diffusée, sans allocation).}, label=lst:helpreq]
// Requête
{
  "attack_id":           "a1b2c3d4-...",
  "helping_peer_id":     "550e8400-...",
  "requesting_node_id":  "node-university"
}

// Réponse 201 Created
{
  "session_id":         "f9e8d7c6-...",
  "attack_id":          "a1b2c3d4-...",
  "helping_peer_id":    "550e8400-...",
  "requesting_node_id": "node-university",
  "status":             "REQUESTED",
  "direction":          "OUTBOUND_REQUEST",
  "requested_at":       "2026-05-22T10:00:05.000Z"
}
\end{lstlisting}

Le pair accepte ou rejette la session après vérification de sa propre capacité disponible~; aucun volume cible ne lui est imposé. S'il accepte, il déclare lui-même le volume qu'il peut offrir (\texttt{accepted\_volume\_gbps})~: c'est cette valeur, et non une capacité issue d'un heartbeat antérieur, qui alimentera le calcul d'allocation WSM.

\begin{lstlisting}[language=json, caption={Acceptation (\texttt{PUT /help/\{id\}/accept}) et rejet (\texttt{PUT /help/\{id\}/reject}).}, label=lst:helpacc]
// Acceptation — corps de la requête
{
  "accepted_volume_gbps": 8.5,
  "tunnel_type":          "GRE",
  "response_time_ms":     230
}
// Réponse 200 OK
{
  "session_id":           "f9e8d7c6-...",
  "status":               "ACCEPTED",
  "accepted_volume_gbps": 8.5,
  "tunnel_type":          "GRE",
  "responded_at":         "2026-05-22T10:00:06.000Z"
}

// Rejet — corps de la requête
{ "rejection_reason": "Capacité insuffisante" }
// Réponse 200 OK
{
  "session_id":       "f9e8d7c6-...",
  "status":           "REJECTED",
  "rejection_reason": "Capacité insuffisante",
  "responded_at":     "2026-05-22T10:00:06.000Z"
}
\end{lstlisting}

% --------------------------------------------------
\subsection{Allocation WSM post-acceptation}
\label{subsec:allocate}

Une fois les réponses des pairs sollicités connues, le nœud victime appelle \texttt{POST /attack/\{id\}/allocate} pour calculer la répartition optimale du flux excédentaire. Ce calcul réutilise exactement l'algorithme WSM (Équation~\ref{eq:wsm}), restreint aux pairs en statut \texttt{ACCEPTED} pour cette attaque, en utilisant leur \texttt{accepted\_volume\_gbps} comme critère $C_p$.

\begin{lstlisting}[language=json, caption={Requête et réponse — \texttt{POST /attack/\{id\}/allocate}.}, label=lst:allocate]
// Requête (aucun corps requis)
POST /api/v1/attack/a1b2c3d4-.../allocate

// Réponse 200 OK
{
  "attack_id": "a1b2c3d4-...",
  "plan": [
    {
      "peer": { "peer_id": "550e8400-...", "peer_name": "node-pme" },
      "wsm_score": 0.612,
      "weight": 0.34,
      "allocation_pct": 34,
      "criteria": { "Cp": 0.71, "Lp_inv": 0.88, "Tp": 0.60, "Rp": 0.50 }
    }
  ]
}
\end{lstlisting}

Le champ \texttt{allocation\_pct} de chaque session \texttt{ACCEPTED} est mis à jour avec le résultat. Cette route peut être rappelée à tout moment pour recalculer le plan sur un ensemble d'accepteurs mis à jour~: c'est notamment le mécanisme réutilisé par la reconsidération automatique de la coalition (section~\ref{subsec:reconsideration}).

% --------------------------------------------------
\subsection{Redirection du trafic et clôture}

Une fois la session acceptée, le nœud victime déclenche la redirection du trafic vers le pair via un tunnel réseau.

\begin{lstlisting}[language=json, caption={Requête et réponse — \texttt{POST /traffic/redirect}.}, label=lst:redirect]
// Requête
{
  "session_id":  "f9e8d7c6-...",
  "tunnel_type": "GRE",
  "volume_gbps": 8.5
}

// Réponse 200 OK
{
  "session_id":        "f9e8d7c6-...",
  "status":            "ACTIVE",
  "tunnel_type":       "GRE",
  "actual_volume_gbps": 8.5,
  "activated_at":      "2026-05-22T10:00:07.000Z"
}
\end{lstlisting}

À la fin de l'attaque, le nœud victime clôture toutes les sessions et déclenche la mise à jour des crédits de réciprocité.

\begin{lstlisting}[language=json, caption={Requête et réponse — \texttt{POST /attack/over}.}, label=lst:over]
// Requête
{
  "attack_id":               "a1b2c3d4-...",
  "session_ids":             ["f9e8d7c6-..."],
  "attack_duration_seconds": 300
}

// Réponse 200 OK
{
  "attack_id":        "a1b2c3d4-...",
  "status":           "ENDED",
  "severity":         "HIGH",
  "ended_at":         "2026-05-22T10:05:00.000Z",
  "duration_seconds": 300,
  "coalition_helped": true,
  "nb_peers_involved": 1
}
\end{lstlisting}

La sévérité finale est calculée selon l'échelle Stormwall~: \texttt{CRITICAL} ($\geq$100 Gbps filtrés), \texttt{HIGH} ($\geq$10 Gbps), \texttt{MEDIUM} ($\geq$1 Gbps), \texttt{LOW} (en dessous).

% --------------------------------------------------
\subsection{Reconsidération automatique de la coalition}
\label{subsec:reconsideration}

Une mitigation en cours n'est jamais figée~: la fonction \texttt{reconsiderCoalition(triggerPeerId, reason)} (\texttt{routes/coalition.js}) est invoquée automatiquement sur trois événements — adhésion d'un nouveau pair, mise à jour de capacité reçue par heartbeat (section~\ref{sec:protocole_com}), ou départ propre (\texttt{POST /peers/goodbye}). Elle ne s'exécute que s'il existe au moins une attaque dont le statut n'est pas \texttt{ENDED}.

\begin{lstlisting}[language=JavaScript, caption={Reconsidération de la coalition sur événement (extrait simplifié).}, label=lst:reconsider]
async function reconsiderCoalition(triggerPeerId, reason) {
  const ongoingAttacks = await Attack.findAll(
    { where: { status: { [Op.ne]: "ENDED" } } });
  if (ongoingAttacks.length === 0) return;

  for (const attack of ongoingAttacks) {
    // Pair tombe pendant une session ACTIVE -> session marquee FAILED
    // Pair toujours actif -> capacite ajustee a la baisse si besoin

    const currentlyCovered = /* somme actual_volume_gbps des sessions ACTIVE */;
    const remainingNeed = attack.overflow_volume_gbps - currentlyCovered;
    if (remainingNeed <= 0) continue;

    // Sollicite de nouveaux candidats (non deja engages) pour combler le manque
    const { plan } = await selectPeers({ overflowGbps: remainingNeed });
    for (const candidate of plan) {
      await solicitPeer(attack, candidate.peer, localNodeId);
    }
  }
}
\end{lstlisting}

Si le pair déclencheur avait une session \texttt{ACTIVE} et qu'il n'est plus joignable (départ ou défaillance), cette session est marquée \texttt{FAILED} et une entrée \texttt{AUDIT\_LOG} est créée. Le manque de couverture résultant (\texttt{overflow\_volume\_gbps} moins le volume encore couvert) déclenche une nouvelle vague de sollicitation (\texttt{solicitPeer}, la même fonction que pour la sollicitation initiale CU~4) auprès des pairs non encore engagés sur cette attaque. Ce mécanisme évite qu'une coalition figée à l'instant $t_0$ devienne obsolète au fil d'une mitigation qui peut durer plusieurs minutes.

% --------------------------------------------------
\section{Tests et évaluation}
\label{sec:tests}

\subsection{Protocole de test}

Les tests se répartissent en deux familles, selon ce qu'elles sollicitent réellement.

\paragraph{Tests d'isolation algorithmique.} \texttt{tests/functional\_test.py} et \texttt{tests/scalability\_test.py} font tourner jusqu'à 500 nœuds sous forme de serveurs \texttt{aiohttp}~\parencite{aiohttp2024} légers, exécutés en mémoire dans un seul processus Python, sans Docker ni PostgreSQL~: chaque nœud simulé réimplémente la forme de l'ensemble des endpoints (enregistrement, heartbeat, sélection WSM, recalcul PeerTrust, sessions d'aide) avec un état purement local. Cette approche isole délibérément la \textbf{complexité du protocole et de l'algorithme} de la variabilité d'une infrastructure réelle (réseau, pool de connexions, E/S disque)~: elle répond à la question \emph{"le protocole et l'algorithme d'allocation scalent-ils en complexité ?"}, pas \emph{"le déploiement Node.js/PostgreSQL tient-il la charge ?"}.

\paragraph{Tests d'intégration réelle.} \texttt{tests/security\_test.py}, \texttt{tests/e2e\_test.py} et \texttt{tests/heartbeat\_test.py} s'exécutent contre le système réellement déployé (\texttt{docker compose up}, Node.js~20~LTS, PostgreSQL~15, sur \texttt{https://localhost:3001})~: tentatives d'accès sans JWT ou avec un token invalide/expiré, parcours complet d'une attaque (sollicitation diffusée, acceptation, allocation WSM, redirection, clôture), reconfiguration automatique de la coalition lors d'une panne de pair, et vérification du heartbeat événementiel sur deux conteneurs réels.

Les deux familles sont complémentaires~: la première caractérise la croissance asymptotique du protocole indépendamment de l'infrastructure, la seconde valide que l'implémentation réelle se comporte conformément à cette caractérisation.

\subsection{Résultats des tests de scalabilité (isolation algorithmique)}

Le Tableau~\ref{tab:scalabilite} présente les résultats collectés via \texttt{scalability\_test.py} (serveurs \texttt{aiohttp} en mémoire, cf.~section précédente)~; la Figure~\ref{fig:scalabilite} en offre une représentation graphique. Ces chiffres caractérisent la complexité du protocole et de l'algorithme d'allocation eux-mêmes, indépendamment de la latence réseau ou des accès disque PostgreSQL~; la démonstration sur le tableau de bord (section~\ref{subsec:demo_dashboard}) valide séparément le comportement du système réellement déployé sur Docker.




\begin{table}[H]
\centering
\caption{Résultats des tests de scalabilité (20 à 500 nœuds).}
\label{tab:scalabilite}
\small
\renewcommand{\arraystretch}{1.35}
\begin{tabular}{|c|c|c|c|c|c|c|c|c|}
\hline
\multirow{2}{*}{\textbf{Nœuds}} &
\multicolumn{3}{c|}{\textbf{Demande d'aide}} &
\multicolumn{3}{c|}{\textbf{Heartbeat}} &
\multicolumn{2}{c|}{\textbf{Allocation du trafic}} \\
\cline{2-9}
& \textbf{OK} & \textbf{Moy (ms)} & \textbf{Max (ms)}
& \textbf{OK} & \textbf{Moy (ms)} & \textbf{Max (ms)}
& \textbf{Durée (ms)} & \textbf{Accept.} \\
\hline
 20 &  20 &   3.42 &   4.26 &  20 &   3.40 &   4.24 & 0.45 &  16 \\
 40 &  40 &   7.61 &   9.26 &  40 &   7.87 &   9.80 & 0.51 &  30 \\
 60 &  60 &  11.55 &  14.73 &  60 &  11.86 &  14.52 & 0.79 &  36 \\
 80 &  80 &  17.22 &  20.50 &  80 &  18.33 &  22.31 & 1.01 &  53 \\
100 & 100 &  21.18 &  25.69 & 100 &  23.19 &  29.21 & 0.82 &  64 \\
150 & 150 &  29.79 &  36.65 & 150 &  29.41 &  34.08 & 0.97 &  99 \\
200 & 200 &  34.76 &  43.50 & 200 &  35.73 &  44.56 & 1.15 & 138 \\
300 & 300 &  71.06 &  98.71 & 300 & 140.91 & 197.70 & 3.30 & 207 \\
500 & 500 & 105.75 & 160.96 & 500 &  98.02 & 152.66 & 3.63 & 335 \\
\hline
\end{tabular}
\end{table}




\begin{figure}[H]
\centering
\figborder[\textwidth]{figures/shieldnet_scalability.png}
\caption{Temps de réponse en fonction du nombre de nœuds simulés (demande d'aide, heartbeat et allocation du trafic).}
\label{fig:scalabilite}
\end{figure}

\subsection{Analyse des résultats}

Quatre observations ressortent des résultats, en gardant à l'esprit qu'il s'agit d'un test d'isolation algorithmique (client et serveurs simulés co-localisés dans un seul processus Python)~: les chiffres caractérisent la forme de la croissance, pas un temps de réponse absolu en production.

\paragraph{Taux de succès de 100\,\% — jusqu'à 500 nœuds.} Aucune erreur n'a été observée sur l'ensemble des neuf paliers testés (20 à 500). Dans ce contexte d'isolation, ce résultat signifie que ni le protocole ni l'algorithme d'allocation n'introduisent de point de rupture structurel à cette échelle~; il ne permet pas de conclure sur la tenue en charge du déploiement réel (Node.js/PostgreSQL/réseau), qui reste à caractériser séparément.

\paragraph{Allocation du trafic~: croissance sous-linéaire.} La durée moyenne de l'allocation (sur 5 répétitions par palier) passe de 0{,}45~ms (16 pairs accepteurs) à 3{,}63~ms (335 pairs accepteurs) — un facteur $\times 8$ pour un nombre de pairs scorés multiplié par $\times 21$. C'est nettement plus favorable que la borne théorique $O(M \log M)$ dérivée pour cet algorithme~: à cette échelle, le calcul lui-même (tri et somme pondérée) reste largement dominé par le coût fixe de la requête, ce qui confirme que l'allocation n'est pas le facteur limitant du protocole.

\paragraph{Demande d'aide et heartbeat~: croissance globalement proportionnelle au nombre de nœuds.} Les deux latences moyennes croissent d'un facteur $\times 31$ et $\times 29$ respectivement entre 20 et 500 nœuds (pour un nombre de nœuds multiplié par $\times 25$) — cohérent avec la complexité $O(N)$ attendue pour $N$ requêtes concurrentes indépendantes, sans dérive quadratique observable. Le palier 300 fait toutefois exception pour le heartbeat (140{,}9~ms, supérieur au palier 500 qui suit) — un rappel que ces mesures, même moyennées sur $N$ requêtes par palier, restent sensibles au bruit d'un environnement d'isolation mono-processus (ordonnancement asyncio, pauses du ramasse-miettes) et doivent être lues comme une tendance, pas comme des valeurs déterministes.

\paragraph{Limite de cette caractérisation.} Ces résultats répondent à la question de la complexité algorithmique, pas à celle de la tenue en charge du système déployé~; les tests d'intégration réelle (section précédente) et la démonstration sur le tableau de bord (section~\ref{subsec:demo_dashboard}) restent la référence pour le comportement du système réellement déployé.

% --------------------------------------------------
\subsection{Résultats des tests de sécurité}

Le Tableau~\ref{tab:auth_tests} présente les cinq scénarios d'attaque simulés et les réponses obtenues du système.









\begin{table}[H]
\centering
\renewcommand{\arraystretch}{1.3}

\begin{tabularx}{\textwidth}{|
>{\raggedright\arraybackslash}X|
>{\centering\arraybackslash}p{2cm}|
>{\centering\arraybackslash}p{1.8cm}|}
\hline
\textbf{Scénario testé} & \textbf{HTTP attendu} & \textbf{Résultat} \\
\hline
Requête sans en-tête \texttt{Authorization} & 401 & \checkmark \\
\hline
Token JWT expiré (émis il y a $>15$~min) & 401 & \checkmark \\
\hline
Token JWT signé avec une clé RSA incorrecte & 401 & \checkmark \\
\hline
Token valide d'un nœud à statut \texttt{BANNED} & 403 & \checkmark \\
\hline
\texttt{peer\_id} inconnu dans le corps de la requête & 404 & \checkmark \\
\hline
\end{tabularx}

\caption{Tests de validation du mécanisme d'authentification}
\label{tab:auth_tests}
\end{table}










Les cinq scénarios ont été rejetés avec le code HTTP correct. Aucun accès non autorisé n'a été accordé~; la vérification séquentielle du middleware (décodage sans vérification → récupération de la clé publique → vérification RS256) a fonctionné dans tous les cas.

% --------------------------------------------------
\subsection{Démonstration sur le tableau de bord de supervision}
\label{subsec:demo_dashboard}

Pour valider l'intégration cohérente de l'ensemble des composants en conditions réalistes, un scénario de bout en bout a été exécuté via le tableau de bord de supervision développé en HTML/JavaScript. Cette interface offre une vision consolidée de l'état de la coalition et permet de déclencher manuellement une simulation d'attaque \gls{ddos} tout en observant en temps réel chaque phase du protocole de mitigation collaborative.

\paragraph{État initial de la coalition.}
La Figure~\ref{fig:cap_overview} présente la vue d'ensemble au moment du lancement de la simulation. La coalition compte 103 pairs enregistrés~: 3~nœuds réels (Datacenter Oran, FAI Algérie Télécom, PME Algéroise) et 100~pairs virtuels répartis en cinq niveaux de confiance. Parmi eux, 93~pairs sont actifs, 23~sont de niveau GOLD et 10~sont bannis. Une attaque précédentes restent ouvertes en base (non clôturées), et 55~sessions d'aide sont encore actives au moment de la capture.

\begin{figure}[H]
\centering
\figborder[\textwidth]{figures/captures/cap_overview.png}
\caption{Vue d'ensemble du tableau de bord~: état de la coalition avant le lancement du scénario.}
\label{fig:cap_overview}
\end{figure}

La Figure~\ref{fig:cap_nodes} présente le tableau des pairs triés par score de confiance décroissant. Les trois nœuds réels (node-datacenter, node-isp, node-pme) apparaissent en tête avec le statut ACTIVE et un score PeerTrust de 1.000 (nœuds GOLD sans historique de défaillance), suivis des pairs virtuels dont les scores varient selon leur niveau (GOLD~: 0.928–0.929, SILVER~: 0.700–0.717).

\begin{figure}[H]
\centering
\figborder[\textwidth]{figures/captures/cap_nodes.png}
\caption{Tableau des pairs de la coalition triés par score PeerTrust décroissant.}
\label{fig:cap_nodes}
\end{figure}

% NOTE : les captures cap_attack_wsm.png / cap_wsm_full.png / cap_mitigation.png
% ont été prises avec l'ancien tableau de bord (sélection WSM avant sollicitation).
% Depuis le passage à la sollicitation diffusée (cf. CU4/CU5bis révisés, chap. 3),
% ces captures doivent être reprises~: l'étape 2 du tableau de bord affiche désormais
% la sollicitation brute (sans score), et l'étape 3 affiche le tableau WSM calculé
% sur les seuls pairs ayant accepté. Les chiffres ci-dessous (38 pairs, 0.905, etc.)
% restent ceux de l'ancienne capture et sont à actualiser après reprise des captures.
\paragraph{Détection et sollicitation de la coalition.}
Une attaque de sévérité HIGH est configurée~: volume de 25~Gbps, protocole UDP, port~80, ciblant le réseau \texttt{192.168.100.0/24}. Avec une capacité locale disponible de 8~Gbps, l'overflow est de 17~Gbps. Le système diffuse automatiquement une demande d'aide (\texttt{POST /help/request}, sans allocation ni volume) à l'ensemble des pairs \texttt{ACTIVE} de la coalition. La Figure~\ref{fig:cap_db_attack_wsm} montre le résultat de cette sollicitation~: une partie des pairs acceptent en déclarant leur propre capacité disponible, et le reste refuse, illustrant le comportement probabiliste implémenté (taux de refus de 15\,\% pour GOLD, 25\,\% pour SILVER, 45\,\% pour BRONZE, 75\,\% pour SUSPECT).

\begin{figure}[H]
\centering
\figborder[\textwidth]{figures/captures/cap_attack_wsm.png}
\caption{Configuration de l'attaque simulée et résultat de la sollicitation diffusée.}
\label{fig:cap_db_attack_wsm}
\end{figure}

\paragraph{Allocation WSM post-acceptation.}
Une fois les réponses connues, le nœud victime appelle \texttt{POST /attack/\{id\}/allocate}~: l'algorithme \gls{wsm} calcule un score composite pour chaque pair ayant accepté, selon la formule $\text{WSM}(p) = 0{,}545 \cdot C + 0{,}193 \cdot L + 0{,}193 \cdot T + 0{,}069 \cdot R$, où $C$ est désormais la capacité que le pair a lui-même déclarée à l'acceptation (\texttt{accepted\_volume\_gbps}), et non plus une valeur issue d'un heartbeat antérieur. La Figure~\ref{fig:cap_wsm_full} présente le classement complet des pairs accepteurs scorés~: les pairs combinant une latence normalisée élevée et une forte capacité déclarée obtiennent les scores les plus élevés et se voient allouer les plus grandes fractions du trafic d'overflow.

\begin{figure}[H]
\centering
\figborder[\textwidth]{figures/captures/cap_wsm_full.png}
\caption{Tableau WSM complet~: classement des pairs accepteurs avec scores et allocations.}
\label{fig:cap_wsm_full}
\end{figure}

\paragraph{Mitigation distribuée active.}
La Figure~\ref{fig:cap_mitigation} présente l'étape de mitigation distribuée. Les tunnels GRE sont activés simultanément pour chaque pair accepteur, selon le plan d'allocation calculé à l'étape précédente, absorbant une fraction substantielle de l'overflow. Les nœuds du tableau sont triés par volume décroissant~: les pairs combinant statut GOLD et forte latence inversée absorbent les volumes les plus importants. À la clôture de l'attaque, chaque pair ayant participé voit son score PeerTrust recalculé avec la sévérité HIGH ($D = 0{,}75$) comme facteur contextuel.

\begin{figure}[H]
\centering
\figborder[\textwidth]{figures/captures/cap_mitigation.png}
\caption{Mitigation distribuée active~: 52 tunnels GRE, 15{,}7~Gbps absorbés, recalcul PeerTrust sur 52 pairs.}
\label{fig:cap_mitigation}
\end{figure}

\paragraph{Traçabilité~: journal d'audit et messages P2P.}
La Figure~\ref{fig:cap_db_audit} présente le journal d'audit horodaté~: les événements \texttt{TRUST\_LEVEL\_CHANGED} confirment l'évolution des niveaux à la clôture de la session — plusieurs pairs progressent de SUSPECT vers BRONZE, et d'autres de BRONZE vers SILVER, témoignant de l'effet de récompense de la participation à une session d'aide réussie.

\begin{figure}[H]
\centering
\figborder[\textwidth]{figures/captures/cap_audit.png}
\caption{Journal d'audit~: événements \texttt{TRUST\_LEVEL\_CHANGED} à la clôture de l'attaque.}
\label{fig:cap_db_audit}
\end{figure}

La Figure~\ref{fig:cap_messages} présente les messages P2P échangés pendant la simulation~: les \texttt{HELP\_ACCEPT} confirment l'acceptation des sessions par les pairs (node-pme, node-isp, node-datacenter, sim-001), et les \texttt{TRAFFIC\_REDIRECT} actent la mise en place des tunnels GRE.

\begin{figure}[H]
\centering
\figborder[\textwidth]{figures/captures/cap_messages.png}
\caption{Messages P2P~: \texttt{HELP\_ACCEPT} et \texttt{TRAFFIC\_REDIRECT} échangés durant la mitigation.}
\label{fig:cap_messages}
\end{figure}

\paragraph{Supervision des heartbeats.}
La Figure~\ref{fig:cap_heartbeats} présente le panel de supervision des heartbeats. Les trois nœuds réels annoncent leur état à node-university sur événement (adhésion à la coalition, puis toute mise à jour de paramètres)~: node-datacenter (charge 20\,\%, 12~Gbps disponibles), node-isp (charge 20\,\%, 16~Gbps disponibles) et node-pme (charge 20\,\%, 4~Gbps disponibles). Ces échanges permettent à chaque nœud de maintenir une vue à jour de la disponibilité réelle de ses pairs, sans qu'aucun minuteur périodique ne génère de trafic en l'absence de changement.

\begin{figure}[H]
\centering
\figborder[\textwidth]{figures/captures/cap_heartbeats.png}
\caption{Panel Heartbeats~: battements périodiques des nœuds réels reçus par node-university.}
\label{fig:cap_heartbeats}
\end{figure}

% --------------------------------------------------
\section{Conclusion}

Ce chapitre a présenté l'implémentation complète du système et ses résultats d'évaluation. Le gestionnaire de confiance implémente fidèlement la formule PeerTrust avec une crédibilité historique par session~: les trois scénarios de validation produisent des scores $T(u)$ conformes aux valeurs théoriques, et le bannissement automatique se déclenche correctement au seuil $T(u) < 0{,}20$. L'algorithme \gls{wsm}, désormais calculé après l'acceptation des pairs sollicités, alloue le flux proportionnellement aux scores composites en s'appuyant sur la capacité que chaque pair déclare lui-même~; un mécanisme de reconsidération automatique (section~\ref{subsec:reconsideration}) recalcule cette allocation chaque fois qu'un pair rejoint, met à jour ses paramètres ou quitte la coalition en cours de mitigation. La double sécurisation mTLS + JWT RS256 a rejeté l'ensemble des cinq scénarios d'attaque simulés avec le code HTTP approprié.

La démonstration sur le tableau de bord de supervision confirme quant à elle l'intégration cohérente de l'ensemble des composants en conditions réalistes~: le scénario complet s'est exécuté avec succès de bout en bout, validant la chaîne authentification~→~détection~→~sollicitation diffusée~→~acceptation~→~allocation \gls{wsm}~→~activation~→~clôture~→~recalcul du score de confiance.
